\documentclass[10pt]{amsart}
\usepackage[a4paper,margin=23mm]{geometry}
\usepackage{amsmath,amssymb,amsthm,mathtools,hyperref,graphicx,xspace,letltxmacro,xpatch}
\usepackage[english]{babel}
\usepackage{xfrac,xcolor}

\usepackage{mathrsfs}
\let\pazocal\mathscr
\multlinegap0pt
\newcommand{\Psfrac}[2]{(\sfrac{#1}{#2})}
\newcommand{\psfrac}[2]{\sfrac{(#1)}{#2}}
\newcommand{\spfrac}[2]{\sfrac{#1}{(#2)}}
\newcommand{\pspfrac}[2]{\sfrac{(#1)}{(#2)}}
\newcommand{\Changel}{\mathcode`l="7160}
\newcommand{\Changelback}{\mathcode`l="716C}
\newcommand{\NoChangel}[1]{%
\expandafter\let\csname old\string#1\endcsname=#1
\let#1=\relax
\newcommand{#1}{\mathcode`l="716C\csname old\string#1\endcsname\mathcode`l="7160 }%
}

\NoChangel{\log}
\NoChangel{\ln}
\NoChangel{\lim}
\NoChangel{\limsup}
\NoChangel{\liminf}
\NoChangel{\varlimsup}
\NoChangel{\varliminf}
\NoChangel{\varinjlim}
\NoChangel{\varprojlim}
\newcommand{\RedefinitSymbole}[1]{%
\expandafter\let\csname old\string#1\endcsname=#1
\let#1=\relax
\newcommand{#1}{\csname old\string#1\endcsname\,}%
}
\RedefinitSymbole{\forall} \RedefinitSymbole{\exists}

\def\mto{\mathchoice{\longmapsto}{\mapsto}{\mapsto}{\mapsto}}
\newcommand{\sbullet}{{\raisebox{1pt}{$\scriptscriptstyle\bullet$}}}
\newcommand{\LtM}{\Lambda^{\!\!\top\!}M}

\usepackage{tikz}
\usetikzlibrary{calc,decorations.markings}

\newtheorem{theorem}{Theorem}[section]
\newtheorem{proposition}[theorem]{Proposition}
\newtheorem{lemma}[theorem]{Lemma}
\newtheorem{corollary}[theorem]{Corollary}

\theoremstyle{definition}
\newtheorem{definition}[theorem]{Definition}
\newtheorem{remark}[theorem]{Remark}
\newtheorem{example}[theorem]{Example}

\newenvironment{refproof}[1]{\begin{proof}[Proof of {#1}]}{\end{proof}}

\newcommand\rr{{\mathbb R}}
\newcommand\zz{{\mathbb Z}}
\newcommand\cc{{\mathbb C}}
\newcommand\nn{{\mathbb N}}

\newcommand\cD{{\pazocal D}}
\newcommand\cU{{\pazocal U}}

\renewcommand\c{\mathrm{c}}
\newcommand{\ph}{\mathrm{cl}}
\newcommand{\GL}{\mathrm{GL}}
\newcommand\loc{\mathrm{loc}}

\let\bra\langle
\let\ket\rangle
\let\diag\Delta
\let\cv\alpha
\let\varalpha s
\let\varm k
\let\asigma a
\let\st\mid
\let\euler X
\let\Feuler V
\let\p\partial
\let\exd u
\let\Im\relax
\let\Re\relax
\DeclareMathOperator{\Im}{Im}
\DeclareMathOperator{\Re}{Re}
\DeclareMathOperator{\wres}{{\scriptstyle W}-res}
\DeclareMathOperator{\resnorm}{res}
\DeclareMathOperator\id{id}
\DeclareMathOperator*{\pf}{fp}
\DeclareMathOperator*{\Vol}{Vol}
\DeclareMathOperator*{\res}{res}

\newcommand{\resdyn}{{\resnorm_\euler}}
\newcommand{\one}{\boldsymbol{1}}
\newcommand*{\defeq}{:=}				
\newcommand\Cf{C^\infty}
\newcommand\zero{{\textit{o}}}
\newcommand\wf{\mathrm{WF}}
\newcommand{\wfl}[2][s]{\mathrm{WF}'^{\,(#1)}_{{\ifempty{#2} h(z) \else \ifnum #2=0 {} \else \ifnum #2=1 \Zweight^{-1} \else \ifnum #2=2 \module{\Im z}^{-1} \else \ifnum #2=3 \bra{\Im z}\ket^{-1} \else \ifnum #2=4 h_1(z) \else \ifnum #2=5 h_2(z) \else \ifnum #2=6 h_1h_2(z) \else \ifnum #2=7 \bra z \ket^{-\infty} \else \ifnum #2=12 \bra z \ket^{-\s12} \fi\fi\fi\fi\fi\fi\fi\fi\fi\fi}}}

\def \ifempty#1{\def\temp{#1} \ifx\temp\empty }

\newcommand{\module}[1]{\left|#1\right|}
\newcommand{\open}[1]{\mathopen{}\mathclose{\left]#1 \right[}}
\newcommand{\clopen}[1]{\mathopen{}\mathclose{\left[#1 \right[}}
\newcommand{\opencl}[1]{\mathopen{}\mathclose{\left]#1 \right]}}
\newcommand{\plancher}[1]{\left \lfloor #1 \right \rfloor }
\renewcommand{\sec}[1]{Section~\ref{#1}}
\newcommand{\secs}[2]{Sections~\ref{#1}--\ref{#2}}
\newcommand{\step}[1]{\subsubsection*{Step #1}}
\renewcommand{\boxed}[1]{#1}

\def\12{\frac{1}{2}}
\def\n2{\frac{n}{2}}
\def\s12{\sfrac{1}{2}}
\def\sn2{\sfrac{n}{2}}
\newcommand\pid{\frac{\pi}{2}}
\newcommand\spid{\sfrac{\pi}{2}}

\newcommand\varvarm{{x_0}}
\newcommand\dvol{\mathop{}\!d\mathrm{vol}}
\newcommand\fantom{\\ &\phantom{=}\,}
\newcommand{\Fs}[2][\cv]{F_{#1}(#2,\vert x\vert_\eta)}
\newcommand{\Fse}[2][\cv]{F_{#1}(#2,\vert{\cdot} \vert_\eta)}
\newcommand{\Fe}[2][\cv]{\mathbf{F}_{#1}(#2, \cdot)}
\newcommand{\Feg}[2][\cv]{\mathbf{F}_{#1}(#2, \vert {\cdot} \vert_g)}

\renewcommand{\t}{{{}^{\mathrm{t}\mkern-2mu}}}

\let\dpl\displaystyle
\providecommand{\skpt}{\par\smallskip\noindent}
\providecommand{\pointir}{.\ }
\providecommand{\eg}{e.g.}\providecommand{\ie}{i.e.}\providecommand{\etc}{etc.}\providecommand{\resp}{resp.}\providecommand{\cf}{cf.}
\providecommand{\mainmatter}{}\providecommand{\backmatter}{}
\let\Subsection\subsection\let\Subsubsection\subsubsection
\setlength{\emergencystretch}{2em}
